\textbf{Phần 1. Trắc nghiệm lựa chọn}
\setcounterpageref{ex}{0}

\begin{ex} %[5H6V32]
Một hình bình hành có độ dài đáy là $24\text{ cm}$, chiều cao tương ứng bằng $\dfrac{3}{4}$ độ dài đáy. Diện tích của hình bình hành đó là
\choice
{$216\text{ cm}^2$}
{$324\text{ cm}^2$}
{\True $432\text{ cm}^2$}
{$864\text{ cm}^2$}
\loigiai{
Chiều cao của hình bình hành là:
$24 \times \dfrac{3}{4} = 18\text{ (cm)}$.\\
Diện tích của hình bình hành là:
$24 \times 18 = 432\text{ (cm}^2\text{)}$.\\
\textbf{Đáp án của câu hỏi là C}
}
\end{ex}

\begin{ex} %[5H6V32]
Một hình thang cân có độ dài hai đáy lần lượt là $14\text{ cm}$ và $26\text{ cm}$, hai cạnh bên mỗi cạnh dài $10\text{ cm}$. Chu vi của hình thang cân đó là
\choice
{$50\text{ cm}$}
{\True $60\text{ cm}$}
{$70\text{ cm}$}
{$80\text{ cm}$}
\loigiai{
Chu vi hình thang cân bằng tổng độ dài 4 cạnh:
$14 + 26 + 10 + 10 = 60\text{ (cm)}$.\\
\textbf{Đáp án của câu hỏi là B}
}
\end{ex}

\begin{ex} %[5H6V32]
Cho hình vuông $ABCD$ có diện tích bằng $64\text{ cm}^2$. Bốn đỉnh của hình vuông nằm trên một đường tròn (đường kính của hình tròn bằng đường chéo hình vuông). Diện tích của hình tròn đó là
\choice
{$64\text{ cm}^2$}
{$80{,}36\text{ cm}^2$}
{\True $100{,}48\text{ cm}^2$}
{$200{,}96\text{ cm}^2$}
\loigiai{
Gọi bán kính hình tròn là $r$. Đường kính hình tròn là $d = 2 \times r$.\\
Diện tích hình vuông tính theo hai đường chéo là: $\dfrac{d \times d}{2} = 64\text{ (cm}^2\text{)}$.\\
Suy ra: $\dfrac{(2 \times r) \times (2 \times r)}{2} = 2 \times r \times r = 64 \implies r \times r = 32$.\\
Diện tích của hình tròn là:
$3{,}14 \times r \times r = 3{,}14 \times 32 = 100{,}48\text{ (cm}^2\text{)}$.\\
\textbf{Đáp án của câu hỏi là C}
}
\end{ex}

\begin{ex} %[5H6V32]
Một mảnh đất gồm một phần hình chữ nhật kích thước $16\text{ m} \times 10\text{ m}$ và một phần hình tam giác vuông tiếp giáp có hai cạnh góc vuông là $6\text{ m}$ và $10\text{ m}$. Diện tích toàn bộ mảnh đất đó là
\choice
{$180\text{ m}^2$}
{\True $190\text{ m}^2$}
{$200\text{ m}^2$}
{$220\text{ m}^2$}
\loigiai{
Diện tích phần hình chữ nhật là:
$16 \times 10 = 160\text{ (m}^2\text{)}$.\\
Diện tích phần hình tam giác vuông là:
$6 \times 10 : 2 = 30\text{ (m}^2\text{)}$.\\
Diện tích toàn bộ mảnh đất là:
$160 + 30 = 190\text{ (m}^2\text{)}$.\\
\textbf{Đáp án của câu hỏi là B}
}
\end{ex}

\begin{ex} %[5H6V32]
Một tam giác có cạnh đáy bằng $20\text{ cm}$. Nếu kéo dài đáy thêm $4\text{ cm}$ thì diện tích tăng thêm $30\text{ cm}^2$. Diện tích tam giác ban đầu là
\choice
{$120\text{ cm}^2$}
{\True $150\text{ cm}^2$}
{$160\text{ cm}^2$}
{$300\text{ cm}^2$}
\loigiai{
Chiều cao của tam giác là:
$30 \times 2 : 4 = 15\text{ (cm)}$.\\
Diện tích tam giác ban đầu là:
$20 \times 15 : 2 = 150\text{ (cm}^2\text{)}$.\\
\textbf{Đáp án của câu hỏi là B}
}
\end{ex}

\begin{ex} %[5H6V32]
Đường kính của một hình tròn bằng $12\text{ cm}$. Chu vi của hình tròn đó là
\choice
{$18{,}84\text{ cm}$}
{\True $37{,}68\text{ cm}$}
{$113{,}04\text{ cm}$}
{$75{,}36\text{ cm}$}
\loigiai{
Chu vi hình tròn là:
$12 \times 3{,}14 = 37{,}68\text{ (cm)}$.\\
\textbf{Đáp án của câu hỏi là B}
}
\end{ex}

\begin{ex} %[5H6V32]
Một hình thoi có hiệu độ dài hai đường chéo là $6\text{ cm}$, tỉ số giữa đường chéo bé và đường chéo lớn là $\dfrac{3}{5}$. Diện tích của hình thoi đó là
\choice
{$135\text{ cm}^2$}
{\True $67{,}5\text{ cm}^2$}
{$90\text{ cm}^2$}
{$45\text{ cm}^2$}
\loigiai{
Hiệu số phần bằng nhau là: $5 - 3 = 2$ (phần).\\
Độ dài đường chéo bé là:
$6 : 2 \times 3 = 9\text{ (cm)}$.\\
Độ dài đường chéo lớn là:
$9 + 6 = 15\text{ (cm)}$.\\
Diện tích hình thoi là:
$9 \times 15 : 2 = 67{,}5\text{ (cm}^2\text{)}$.\\
\textbf{Đáp án của câu hỏi là B}
}
\end{ex}

\begin{ex} %[5H6V32]
Cho hình thang có đáy lớn $18\text{ cm}$, đáy bé $12\text{ cm}$. Chiều cao bằng trung bình cộng hai đáy. Diện tích hình thang là
\choice
{$150\text{ cm}^2$}
{\True $225\text{ cm}^2$}
{$450\text{ cm}^2$}
{$180\text{ cm}^2$}
\loigiai{
Chiều cao hình thang là:
$(18 + 12) : 2 = 15\text{ (cm)}$.\\
Diện tích hình thang là:
$(18 + 12) \times 15 : 2 = 225\text{ (cm}^2\text{)}$.\\
\textbf{Đáp án của câu hỏi là B}
}
\end{ex}

\begin{ex} %[5H6V32]
Nếu tăng bán kính của một hình tròn lên gấp 3 lần thì chu vi của hình tròn đó
\choice
{Tăng gấp 9 lần}
{Tăng gấp 6 lần}
{\True Tăng gấp 3 lần}
{Không thay đổi}
\loigiai{
Chu vi hình tròn $C = 2 \times 3{,}14 \times r$, tỉ lệ thuận với bán kính $r$.\\
Do đó khi bán kính tăng gấp 3 lần thì chu vi cũng tăng gấp 3 lần.\\
\textbf{Đáp án của câu hỏi là C}
}
\end{ex}

\begin{ex} %[5H6V32]
Một hình chữ nhật có chu vi $50\text{ cm}$, chiều dài hơn chiều rộng $7\text{ cm}$. Diện tích của hình chữ nhật đó là
\choice
{$120\text{ cm}^2$}
{\True $144\text{ cm}^2$}
{$150\text{ cm}^2$}
{$168\text{ cm}^2$}
\loigiai{
Nửa chu vi hình chữ nhật là:
$50 : 2 = 25\text{ (cm)}$.\\
Chiều rộng hình chữ nhật là:
$(25 - 7) : 2 = 9\text{ (cm)}$.\\
Chiều dài hình chữ nhật là:
$25 - 9 = 16\text{ (cm)}$.\\
Diện tích hình chữ nhật là:
$16 \times 9 = 144\text{ (cm}^2\text{)}$.\\
\textbf{Đáp án của câu hỏi là B}
}
\end{ex}

\begin{ex} %[5H6V32]
Trong các hình phẳng sau: hình tam giác đều, hình chữ nhật, hình thoi, hình vuông; hình nào có đúng $4$ trục đối xứng?
\choice
{Hình tam giác đều}
{Hình chữ nhật}
{Hình thoi}
{\True Hình vuông}
\loigiai{
Hình vuông có 4 trục đối xứng (2 đường trung trực của các cạnh đối và 2 đường chéo).\\
\textbf{Đáp án của câu hỏi là D}
}
\end{ex}

\begin{ex} %[5H6V32]
Một biển báo giao thông hình tròn có bán kính bằng $25\text{ cm}$. Diện tích của biển báo đó là
\choice
{$157\text{ cm}^2$}
{$78{,}5\text{ cm}^2$}
{\True $1\,962{,}5\text{ cm}^2$}
{$3\,925\text{ cm}^2$}
\loigiai{
Diện tích biển báo là:
$3{,}14 \times 25 \times 25 = 1\,962{,}5\text{ (cm}^2\text{)}$.\\
\textbf{Đáp án của câu hỏi là C}
}
\end{ex}

\textbf{Phần 2. Trắc nghiệm đúng sai}
\setcounterpageref{ex}{0}

\begin{ex} %[5H6V32]
Cho tam giác $MNP$ vuông tại $M$ có độ dài hai cạnh góc vuông $MN = 6\text{ cm}, MP = 8\text{ cm}$ và cạnh huyền $NP = 10\text{ cm}$.
\choiceTF
{\True Diện tích của tam giác $MNP$ bằng $24\text{ cm}^2$}
{\True Chu vi của tam giác $MNP$ bằng $24\text{ cm}$}
{Chiều cao kẻ từ đỉnh $M$ vuông góc với cạnh huyền $NP$ dài đúng $5\text{ cm}$}
{\True Diện tích tam giác $MNP$ bằng một nửa diện tích hình chữ nhật có kích thước $6\text{ cm} \times 8\text{ cm}$}
\loigiai{
\begin{itemchoice}
\item Đúng. Diện tích: $6 \times 8 : 2 = 24\text{ (cm}^2\text{)}$.
\item Đúng. Chu vi: $6 + 8 + 10 = 24\text{ (cm)}$.
\item Sai. Chiều cao kẻ từ $M$ là: $24 \times 2 : 10 = 4{,}8\text{ (cm)} 
eq 5\text{ cm}$.
\item Đúng. Diện tích hình chữ nhật là $6 \times 8 = 48\text{ cm}^2$, tam giác có diện tích $24\text{ cm}^2$.
\end{itemchoice}
}
\end{ex}

\begin{ex} %[5H6V32]
Cho hình thang vuông $ABCD$ có các góc ở đỉnh $A$ và $D$ là góc vuông, đáy bé $AB = 8\text{ cm}$, đáy lớn $CD = 14\text{ cm}$ và cạnh bên $AD = 6\text{ cm}$.
\choiceTF
{\True Diện tích của hình thang $ABCD$ bằng $66\text{ cm}^2$}
{\True Cạnh bên $AD$ vừa là cạnh bên vừa là chiều cao của hình thang}
{\True Hạ đường cao $BH$ vuông góc với $CD$, tứ giác $ABHD$ là hình chữ nhật có diện tích $48\text{ cm}^2$}
{Hình tam giác vuông $BHC$ có diện tích bằng $20\text{ cm}^2$}
\loigiai{
\begin{itemchoice}
\item Đúng. Diện tích: $(8 + 14) \times 6 : 2 = 66\text{ (cm}^2\text{)}$.
\item Đúng. Vì $AD \perp AB$ và $AD \perp CD$ nên $AD$ là chiều cao hình thang.
\item Đúng. Tứ giác $ABHD$ có 4 góc vuông nên là hình chữ nhật kích thước $8\text{ cm} \times 6\text{ cm}$, diện tích là $48\text{ cm}^2$.
\item Sai. Đáy $HC = 14 - 8 = 6\text{ (cm)}$, diện tích tam giác $BHC$ là $6 \times 6 : 2 = 18\text{ (cm}^2\text{)} 
eq 20\text{ cm}^2$.
\end{itemchoice}
}
\end{ex}

\textbf{Phần 3. Trắc nghiệm trả lời ngắn}
\setcounterpageref{ex}{0}

\begin{ex} %[5H6V32]
Một hình bình hành có diện tích bằng $180\text{ cm}^2$ và chiều cao là $12\text{ cm}$. Độ dài cạnh đáy tương ứng của hình bình hành đó là bao nhiêu xăng-ti-mét?
\shortans{15}
\loigiai{
Độ dài cạnh đáy hình bình hành là:
$180 : 12 = 15\text{ (cm)}$.
}
\end{ex}

\begin{ex} %[5H6V32]
Một hình thoi có diện tích bằng $96\text{ cm}^2$, một đường chéo dài $12\text{ cm}$. Đường chéo còn lại dài bao nhiêu xăng-ti-mét?
\shortans{16}
\loigiai{
Độ dài đường chéo còn lại là:
$96 \times 2 : 12 = 16\text{ (cm)}$.
}
\end{ex}

\begin{ex} %[5H6V32]
Một hình tròn có diện tích bằng $78{,}5\text{ cm}^2$. Tính đường kính của hình tròn đó theo đơn vị xăng-ti-mét.
\shortans{10}
\loigiai{
Tích hai bán kính là:
$r \times r = 78{,}5 : 3{,}14 = 25$.\\
Vì $5 \times 5 = 25$ nên bán kính hình tròn là $5\text{ cm}$.\\
Đường kính hình tròn là:
$5 \times 2 = 10\text{ (cm)}$.
}
\end{ex}

\begin{ex} %[5H6V32]
Một hình tam giác có cạnh đáy dài $24\text{ cm}$ và chiều cao tương ứng bằng $15\text{ cm}$. Diện tích của hình tam giác đó là bao nhiêu xăng-ti-mét vuông?
\shortans{180}
\loigiai{
Diện tích hình tam giác là:
$24 \times 15 : 2 = 180\text{ (cm}^2\text{)}$.
}
\end{ex}

\begin{ex} %[5H6V32]
Một hình thang có đáy lớn $25\text{ m}$, đáy bé $15\text{ m}$ và chiều cao $10\text{ m}$. Diện tích hình thang đó bằng bao nhiêu mét vuông?
\shortans{200}
\loigiai{
Diện tích hình thang là:
$(25 + 15) \times 10 : 2 = 200\text{ (m}^2\text{)}$.
}
\end{ex}

\begin{ex} %[5H6V32]
Cho một hình chữ nhật có chiều dài $20\text{ cm}$ và chiều rộng $14\text{ cm}$. Người ta vẽ một hình thoi có các đỉnh là trung điểm bốn cạnh của hình chữ nhật đó. Tính diện tích của hình thoi theo đơn vị xăng-ti-mét vuông.
\shortans{140}
\loigiai{
Hai đường chéo của hình thoi bằng chiều dài và chiều rộng hình chữ nhật ($20\text{ cm}$ và $14\text{ cm}$).\\
Diện tích của hình thoi là:
$20 \times 14 : 2 = 140\text{ (cm}^2\text{)}$.
}
\end{ex}
